% =============================================================================
%  Fundamentos compartilhados --- Módulo 14: teste de software
%  Incluído pelas trilhas Kotlin e Flutter com % =============================================================================
%  Fundamentos compartilhados --- Módulo 14: teste de software
%  Incluído pelas trilhas Kotlin e Flutter com % =============================================================================
%  Fundamentos compartilhados --- Módulo 14: teste de software
%  Incluído pelas trilhas Kotlin e Flutter com % =============================================================================
%  Fundamentos compartilhados --- Módulo 14: teste de software
%  Incluído pelas trilhas Kotlin e Flutter com \input{../fundamentos/f14-testes}
% =============================================================================
\section{Fundamentos de teste de software}

\begin{frame}{Erro, defeito e falha}
  \begin{definicao}[Três palavras que não são sinônimos]
    \textbf{Erro} (\emph{error}): engano cometido por uma pessoa.
    \textbf{Defeito} (\emph{defect}, \emph{bug}): imperfeição no código ou em outro artefato, consequência do erro.
    \textbf{Falha} (\emph{failure}): comportamento incorreto observado quando o defeito é executado.
  \end{definicao}
  \centering
  \begin{tikzpicture}[node distance=17mm]
    \node[caixa,minimum width=3.0cm] (e) {\textbf{Erro}\\ \tiny o dev usa \texttt{>} em vez de \texttt{>=}};
    \node[caixapetroleo,minimum width=3.0cm,right=of e] (d) {\textbf{Defeito}\\ \tiny a regra de alerta fica errada};
    \node[caixaambar,minimum width=3.2cm,right=of d] (f) {\textbf{Falha}\\ \tiny a 30,0\,°C nenhum alerta aparece};
    \draw[seta] (e) -- node[rotulo,above]{\tiny introduz} (d);
    \draw[seta] (d) -- node[rotulo,above]{\tiny executado} (f);
  \end{tikzpicture}
  \vspace{2pt}
  \begin{columns}[T]
    \begin{column}{0.48\textwidth}
      \begin{conceito}[Consequência]
        Um defeito pode ficar anos sem causar falha, se a linha nunca for executada nas condições certas. Testar é \textbf{provocar} essas condições.
      \end{conceito}
    \end{column}
    \begin{column}{0.48\textwidth}
      \begin{historia}
        Em 1947, a equipe do computador Harvard Mark II, da qual Grace Hopper fazia parte, achou uma mariposa presa num relé e a colou no diário de bordo. A palavra \emph{bug} já existia, mas o episódio a popularizou.
      \end{historia}
    \end{column}
  \end{columns}
\end{frame}

\begin{frame}{Para que serve testar}
  \begin{columns}[T]
    \begin{column}{0.5\textwidth}
      \begin{conceito}[Objetivos dos testes]
        \begin{itemize}
          \item Encontrar defeitos antes do usuário.
          \item Dar confiança para entregar.
          \item Impedir que defeitos corrigidos voltem (\textbf{regressão}).
          \item Documentar, em código executável, como o sistema deve se comportar.
          \item Permitir \textbf{refatorar} sem medo.
        \end{itemize}
      \end{conceito}
      \begin{atencao}[Princípio fundamental]
        Testes mostram que defeitos \textbf{existem}; nunca provam que \textbf{não existem}. A ideia é de Edsger Dijkstra.
      \end{atencao}
    \end{column}
    \begin{column}{0.48\textwidth}
      \begin{pordentro}[Teste exaustivo é impossível]
        Uma função com dois inteiros de 32 bits tem $2^{64}$ combinações de entrada, cerca de $1{,}8\times10^{19}$. A um bilhão de testes por segundo, seriam quase 600 anos. Por isso escolhemos casos \textbf{representativos} com técnicas de projeto de testes.
      \end{pordentro}
      \begin{analogia}
        O \emph{crash test} de um carro não bate o veículo em todas as velocidades e ângulos: escolhe poucos cenários que representam os acidentes reais.
      \end{analogia}
    \end{column}
  \end{columns}
\end{frame}

\begin{frame}{O custo de achar tarde}
  \begin{columns}[c]
    \begin{column}{0.52\textwidth}
      \centering
      \begin{tikzpicture}[x=0.95cm,y=0.5cm]
        \draw[cinzamedio] (0,0) -- (5.3,0);
        \foreach \x/\h/\t in {0.4/0.6/Requisitos,1.4/1.2/Projeto,2.4/2.2/Código,3.4/3.8/Teste,4.4/6.4/Produção}{
          \fill[azulpetroleo] (\x,0) rectangle ++(0.6,\h);
          \node[font=\tiny,rotate=35,anchor=north east] at (\x+0.45,-0.1) {\t};}
        \fill[ambar] (4.4,0) rectangle ++(0.6,6.4);
        \node[rotulo,font=\tiny,align=center] at (2.2,6.5) {custo relativo de corrigir\\ (ilustrativo)};
      \end{tikzpicture}
    \end{column}
    \begin{column}{0.46\textwidth}
      \begin{conceito}[Por que custa mais]
        Um defeito achado em produção exige investigar relatos, reproduzir, corrigir, testar de novo, gerar versão, passar pela revisão da loja e esperar os usuários atualizarem. No app, ainda há dados já gravados errado.
      \end{conceito}
      \begin{dica}[\emph{Shift left}]
        Antecipar a qualidade: revisar requisitos, escrever testes junto com o código e rodá-los a cada \emph{commit}.
      \end{dica}
      \begin{historia}
        Barry Boehm mediu esse efeito nos anos 1970 e 1980; os números variam entre projetos, mas a tendência se repete.
      \end{historia}
    \end{column}
  \end{columns}
\end{frame}

\begin{frame}{Verificação, validação e níveis de teste}
  \begin{columns}[c]
    \begin{column}{0.56\textwidth}
      \centering
      \begin{tikzpicture}[every node/.style={font=\scriptsize,minimum width=2.25cm,minimum height=6mm}]
        \node[caixa] (l1) at (1.15,0) {Requisitos};
        \node[caixa] (l2) at (1.55,-0.75) {Especificação};
        \node[caixa] (l3) at (1.95,-1.5) {Arquitetura};
        \node[caixa] (l4) at (2.35,-2.25) {Projeto detalhado};
        \node[caixaescura] (c) at (3.8,-3.0) {Código};
        \node[caixaambar] (r4) at (5.25,-2.25) {Teste de unidade};
        \node[caixaambar] (r3) at (5.65,-1.5) {Integração};
        \node[caixaambar] (r2) at (6.05,-0.75) {Teste de sistema};
        \node[caixaambar] (r1) at (6.45,0) {Aceitação};
        \foreach \a/\b in {l1/l2,l2/l3,l3/l4,l4/c,c/r4,r4/r3,r3/r2,r2/r1} \draw[seta] (\a) -- (\b);
        \foreach \a/\b in {l1/r1,l2/r2,l3/r3,l4/r4} \draw[pontilhado,<->] (\a) -- (\b);
      \end{tikzpicture}\\[2pt]
      {\tiny Modelo em V: cada nível de teste confere um nível de especificação.}
    \end{column}
    \begin{column}{0.42\textwidth}
      \begin{definicao}[Verificação]
        ``Estamos construindo o produto \textbf{corretamente}?'' Confere o software contra a especificação.
      \end{definicao}
      \begin{definicao}[Validação]
        ``Estamos construindo o produto \textbf{certo}?'' Confere se o software resolve o problema real do usuário.
      \end{definicao}
      \begin{dica}
        O teste de aceitação é feito com o cliente: no LabMonitor, os técnicos do laboratório usando o app.
      \end{dica}
    \end{column}
  \end{columns}
\end{frame}

\begin{frame}{Caixa-preta: partições e valores-limite}
  \begin{columns}[T]
    \begin{column}{0.5\textwidth}
      \begin{definicao}[Teste caixa-preta]
        Projetado a partir da \textbf{especificação}, sem olhar o código: entradas e saídas esperadas.
      \end{definicao}
      \begin{conceito}[Partição de equivalência]
        Divide as entradas em classes que o sistema deve tratar do mesmo jeito; basta um representante de cada classe.
      \end{conceito}
      \begin{conceito}[Análise de valor-limite]
        Defeitos se concentram nas bordas (\texttt{<} no lugar de \texttt{<=}). Teste o limite e seus vizinhos imediatos.
      \end{conceito}
    \end{column}
    \begin{column}{0.48\textwidth}
      \begin{exemplo}[Regra: a reserva dura de 1 a 4 horas inteiras]
        \centering
        \begin{tikzpicture}[x=0.62cm]
          \draw[cinzamedio,-{Stealth}] (-0.5,0) -- (6.5,0);
          \fill[vermelhoerro!25] (-0.5,-0.15) rectangle (0.5,0.15);
          \fill[verdeok!35] (0.5,-0.15) rectangle (4.5,0.15);
          \fill[vermelhoerro!25] (4.5,-0.15) rectangle (6.5,0.15);
          \foreach \v in {0,1,2,3,4,5,6} \node[font=\tiny,below] at (\v,-0.15) {\v};
          \node[font=\tiny,above] at (0,0.15) {inválida};
          \node[font=\tiny,above] at (2.5,0.15) {válida};
          \node[font=\tiny,above] at (5.5,0.15) {inválida};
        \end{tikzpicture}\par
        \scriptsize
        \begin{tabular}{@{}l l@{}}
          \toprule
          \cabecalho{Casos} & \cabecalho{Por quê} \\
          \midrule
          2 h & representante da classe válida \\
          0 h e 1 h & limite inferior e vizinho \\
          4 h e 5 h & limite superior e vizinho \\
          $-1$ h, ``abc'' & entradas absurdas \\
          \bottomrule
        \end{tabular}
      \end{exemplo}
    \end{column}
  \end{columns}
\end{frame}

\begin{frame}[fragile]{Caixa-branca: critérios de cobertura}
  \begin{columns}[T]
    \begin{column}{0.46\textwidth}
      \begin{definicao}[Teste caixa-branca]
        Projetado a partir da \textbf{estrutura interna} do código: comandos, decisões e caminhos.
      \end{definicao}
\begin{codigo}[title=Pseudocódigo]{bash}
se temp >= limite E alertaAtivo:
    notificar()
registrar(temp)
\end{codigo}
      \begin{atencao}
        Cobertura mede o que foi \textbf{executado}, não o que foi \textbf{conferido}: um teste sem \cod{expect} também ``cobre''.
      \end{atencao}
    \end{column}
    \begin{column}{0.52\textwidth}
      \scriptsize
      \begin{tabularx}{\linewidth}{@{}p{2.1cm} L@{}}
        \toprule
        \cabecalho{Critério} & \cabecalho{Exige, no exemplo} \\
        \midrule
        \textbf{Comandos} & cada linha rodar ao menos uma vez: 1 teste (30\,°C, ativo) basta \\
        \textbf{Decisões} & cada \cod{se} (\emph{branch}) ser verdadeiro e falso: 2 testes \\
        \textbf{Condições} & cada condição simples verdadeira e falsa: \cod{temp} e \cod{alertaAtivo} variando \\
        \textbf{Caminhos} & todas as rotas possíveis; cresce exponencialmente com laços e \cod{se}s \\
        \bottomrule
      \end{tabularx}
      \begin{conceito}[Hierarquia]
        Caminhos $\Rightarrow$ decisões $\Rightarrow$ comandos: quem satisfaz um critério mais forte satisfaz os mais fracos.
      \end{conceito}
    \end{column}
  \end{columns}
\end{frame}

\begin{frame}{Dublês de teste}
  \begin{columns}[T]
    \begin{column}{0.62\textwidth}
      \scriptsize
      \begin{tabularx}{\linewidth}{@{}l L L@{}}
        \toprule
        \cabecalho{Dublê} & \cabecalho{O que é} & \cabecalho{Exemplo no LabMonitor} \\
        \midrule
        \textbf{Dummy} & objeto só para preencher um parâmetro; nunca é usado & um \emph{logger} vazio exigido pelo construtor \\
        \textbf{Stub} & devolve respostas prontas & repositório que sempre devolve 3 bancadas \\
        \textbf{Spy} & stub que também registra as chamadas & conta quantas notificações foram enviadas \\
        \textbf{Mock} & objeto com \textbf{expectativas} de interação, verificadas no teste & ``\cod{reservar(7)} deve ser chamado uma vez'' \\
        \textbf{Fake} & implementação funcional simplificada & banco em memória, API falsa com dados fixos \\
        \bottomrule
      \end{tabularx}
    \end{column}
    \begin{column}{0.36\textwidth}
      \begin{definicao}[Dublê de teste]
        Objeto que substitui uma dependência real (rede, banco, sensor, relógio) durante o teste.
      \end{definicao}
      \begin{analogia}
        Como o dublê de cinema: entra nas cenas perigosas ou caras no lugar do ator.
      \end{analogia}
      \begin{historia}
        A classificação é de Gerard Meszaros, no livro \emph{xUnit Test Patterns} (2007).
      \end{historia}
    \end{column}
  \end{columns}
\end{frame}

\begin{frame}{Propriedades de um bom teste}
  \begin{columns}[T]
    \begin{column}{0.5\textwidth}
      \scriptsize
      \begin{tabularx}{\linewidth}{@{}l L@{}}
        \toprule
        \cabecalho{FIRST} & \cabecalho{Significado} \\
        \midrule
        \textbf{F}ast & rápido: roda em milissegundos, então roda sempre \\
        \textbf{I}ndependent & não depende da ordem nem de outros testes \\
        \textbf{R}epeatable & mesmo resultado em qualquer máquina e horário \\
        \textbf{S}elf-validating & passa ou falha sozinho, sem alguém olhar a saída \\
        \textbf{T}imely & escrito junto com o código (ou antes) \\
        \bottomrule
      \end{tabularx}
      \begin{conceito}[Estrutura]
        \textbf{AAA}: \emph{Arrange}, \emph{Act}, \emph{Assert}.\\
        \textbf{BDD}: \emph{Given}, \emph{When}, \emph{Then} (dado, quando, então).
      \end{conceito}
    \end{column}
    \begin{column}{0.48\textwidth}
      \begin{erro}[Teste instável (\emph{flaky})]
        Passa e falha sem o código mudar. Causas típicas: depender do relógio real, da rede, da ordem de execução, de esperas fixas (\emph{sleep}) ou de estado compartilhado entre testes.
      \end{erro}
      \begin{dica}[Teste também é código]
        Merece nomes claros, nada de duplicação e revisão. Um bom nome de teste descreve uma \textbf{regra}: ``reserva sobreposta é recusada''.
      \end{dica}
    \end{column}
  \end{columns}
\end{frame}

\begin{frame}[fragile]{TDD, BDD e regressão}
  \begin{columns}[T]
    \begin{column}{0.4\textwidth}
      \centering
      \begin{tikzpicture}
        \node[circle,fill=vermelhoerro,text=white,font=\scriptsize\bfseries,minimum size=1.35cm,align=center] (v) at (90:1.35) {Vermelho\\ \tiny teste falha};
        \node[circle,fill=verdeok,text=white,font=\scriptsize\bfseries,minimum size=1.35cm,align=center] (g) at (-30:1.35) {Verde\\ \tiny passa};
        \node[circle,fill=azulpetroleo,text=white,font=\scriptsize\bfseries,minimum size=1.35cm,align=center] (r) at (210:1.35) {Refatorar\\ \tiny melhora};
        \draw[setaambar] (v) to[bend left=25] (g);
        \draw[setaambar] (g) to[bend left=25] (r);
        \draw[setaambar] (r) to[bend left=25] (v);
      \end{tikzpicture}
      \begin{definicao}[TDD]
        \emph{Test-Driven Development} (Kent Beck): o teste é escrito \textbf{antes} do código e guia o projeto.
      \end{definicao}
    \end{column}
    \begin{column}{0.58\textwidth}
\begin{codigo}[title=BDD: cenário em linguagem natural (Gherkin)]{bash}
Funcionalidade: Reserva de bancada
  Cenário: horário já ocupado
    Dado que a bancada B7 está reservada das 14h às 16h
    Quando eu peço a B7 das 15h às 17h
    Então vejo "Horário indisponível"
\end{codigo}
      \begin{definicao}[BDD]
        \emph{Behavior-Driven Development}: descreve o comportamento em exemplos que clientes, testadores e devs entendem; os cenários viram testes.
      \end{definicao}
      \begin{definicao}[Teste de regressão]
        Rodar de novo os testes existentes após cada mudança, para garantir que nada que funcionava quebrou. Automatizar é o que torna isso viável.
      \end{definicao}
    \end{column}
  \end{columns}
\end{frame}

\begin{frame}{Desafios do teste em apps móveis}
  \begin{columns}[T]
    \begin{column}{0.5\textwidth}
      \scriptsize
      \begin{tabularx}{\linewidth}{@{}p{2.0cm} L@{}}
        \toprule
        \cabecalho{Desafio} & \cabecalho{O que pode dar errado} \\
        \midrule
        Fragmentação & milhares de modelos Android, versões de sistema, tamanhos e densidades de tela \\
        Rede & sem sinal, 3G lento, troca Wi-Fi/móvel no meio de uma requisição \\
        Interrupções & ligação, notificação, app enviado ao fundo e encerrado pelo sistema \\
        Hardware & câmera, GPS, Bluetooth e sensores diferentes ou ausentes \\
        Permissões & usuário nega, concede ``só desta vez'' ou revoga depois \\
        Recursos & bateria fraca, pouca memória, armazenamento cheio \\
        Pessoas & idioma, fonte grande, leitor de tela, modo escuro \\
        \bottomrule
      \end{tabularx}
    \end{column}
    \begin{column}{0.48\textwidth}
      \begin{conceito}[Estratégia]
        \begin{itemize}
          \item Automatizar a lógica em testes de unidade e widget (rápidos, sem aparelho).
          \item Definir uma \textbf{matriz de dispositivos} representativa (versões antigas e novas, tela pequena e grande).
          \item Usar emuladores no dia a dia e aparelhos reais ou \emph{device farms} antes de publicar.
          \item Simular rede ruim e interrupções.
        \end{itemize}
      \end{conceito}
      \begin{dica}
        As estatísticas do Android Studio e do console da loja mostram quais versões seus usuários realmente usam.
      \end{dica}
    \end{column}
  \end{columns}
\end{frame}

\begin{frame}{Qualidade é mais do que testar}
  \begin{columns}[T]
    \begin{column}{0.5\textwidth}
      \begin{definicao}[Modelo de qualidade ISO/IEC 25010]
        Organiza a qualidade de um produto de software em características como adequação funcional, eficiência de desempenho, compatibilidade, usabilidade, confiabilidade, segurança, manutenibilidade e portabilidade.
      \end{definicao}
      \begin{dica}
        A revisão de 2023 renomeou algumas características (usabilidade virou capacidade de interação) e incluiu \emph{safety}.
      \end{dica}
    \end{column}
    \begin{column}{0.48\textwidth}
      \begin{conceito}[Técnicas que se somam]
        \begin{itemize}
          \item \textbf{Estáticas} (sem executar): revisão de código, análise estática, \emph{lint}, tipagem forte.
          \item \textbf{Dinâmicas} (executando): testes automatizados e exploratórios.
          \item \textbf{Em produção}: monitoramento de falhas, métricas e \emph{feedback} dos usuários.
        \end{itemize}
      \end{conceito}
      \begin{destaque}
        Qualidade se \textbf{constrói} durante o desenvolvimento; testes confirmam que ela está lá.
      \end{destaque}
    \end{column}
  \end{columns}
\end{frame}

% =============================================================================
\section{Fundamentos de teste de software}

\begin{frame}{Erro, defeito e falha}
  \begin{definicao}[Três palavras que não são sinônimos]
    \textbf{Erro} (\emph{error}): engano cometido por uma pessoa.
    \textbf{Defeito} (\emph{defect}, \emph{bug}): imperfeição no código ou em outro artefato, consequência do erro.
    \textbf{Falha} (\emph{failure}): comportamento incorreto observado quando o defeito é executado.
  \end{definicao}
  \centering
  \begin{tikzpicture}[node distance=17mm]
    \node[caixa,minimum width=3.0cm] (e) {\textbf{Erro}\\ \tiny o dev usa \texttt{>} em vez de \texttt{>=}};
    \node[caixapetroleo,minimum width=3.0cm,right=of e] (d) {\textbf{Defeito}\\ \tiny a regra de alerta fica errada};
    \node[caixaambar,minimum width=3.2cm,right=of d] (f) {\textbf{Falha}\\ \tiny a 30,0\,°C nenhum alerta aparece};
    \draw[seta] (e) -- node[rotulo,above]{\tiny introduz} (d);
    \draw[seta] (d) -- node[rotulo,above]{\tiny executado} (f);
  \end{tikzpicture}
  \vspace{2pt}
  \begin{columns}[T]
    \begin{column}{0.48\textwidth}
      \begin{conceito}[Consequência]
        Um defeito pode ficar anos sem causar falha, se a linha nunca for executada nas condições certas. Testar é \textbf{provocar} essas condições.
      \end{conceito}
    \end{column}
    \begin{column}{0.48\textwidth}
      \begin{historia}
        Em 1947, a equipe do computador Harvard Mark II, da qual Grace Hopper fazia parte, achou uma mariposa presa num relé e a colou no diário de bordo. A palavra \emph{bug} já existia, mas o episódio a popularizou.
      \end{historia}
    \end{column}
  \end{columns}
\end{frame}

\begin{frame}{Para que serve testar}
  \begin{columns}[T]
    \begin{column}{0.5\textwidth}
      \begin{conceito}[Objetivos dos testes]
        \begin{itemize}
          \item Encontrar defeitos antes do usuário.
          \item Dar confiança para entregar.
          \item Impedir que defeitos corrigidos voltem (\textbf{regressão}).
          \item Documentar, em código executável, como o sistema deve se comportar.
          \item Permitir \textbf{refatorar} sem medo.
        \end{itemize}
      \end{conceito}
      \begin{atencao}[Princípio fundamental]
        Testes mostram que defeitos \textbf{existem}; nunca provam que \textbf{não existem}. A ideia é de Edsger Dijkstra.
      \end{atencao}
    \end{column}
    \begin{column}{0.48\textwidth}
      \begin{pordentro}[Teste exaustivo é impossível]
        Uma função com dois inteiros de 32 bits tem $2^{64}$ combinações de entrada, cerca de $1{,}8\times10^{19}$. A um bilhão de testes por segundo, seriam quase 600 anos. Por isso escolhemos casos \textbf{representativos} com técnicas de projeto de testes.
      \end{pordentro}
      \begin{analogia}
        O \emph{crash test} de um carro não bate o veículo em todas as velocidades e ângulos: escolhe poucos cenários que representam os acidentes reais.
      \end{analogia}
    \end{column}
  \end{columns}
\end{frame}

\begin{frame}{O custo de achar tarde}
  \begin{columns}[c]
    \begin{column}{0.52\textwidth}
      \centering
      \begin{tikzpicture}[x=0.95cm,y=0.5cm]
        \draw[cinzamedio] (0,0) -- (5.3,0);
        \foreach \x/\h/\t in {0.4/0.6/Requisitos,1.4/1.2/Projeto,2.4/2.2/Código,3.4/3.8/Teste,4.4/6.4/Produção}{
          \fill[azulpetroleo] (\x,0) rectangle ++(0.6,\h);
          \node[font=\tiny,rotate=35,anchor=north east] at (\x+0.45,-0.1) {\t};}
        \fill[ambar] (4.4,0) rectangle ++(0.6,6.4);
        \node[rotulo,font=\tiny,align=center] at (2.2,6.5) {custo relativo de corrigir\\ (ilustrativo)};
      \end{tikzpicture}
    \end{column}
    \begin{column}{0.46\textwidth}
      \begin{conceito}[Por que custa mais]
        Um defeito achado em produção exige investigar relatos, reproduzir, corrigir, testar de novo, gerar versão, passar pela revisão da loja e esperar os usuários atualizarem. No app, ainda há dados já gravados errado.
      \end{conceito}
      \begin{dica}[\emph{Shift left}]
        Antecipar a qualidade: revisar requisitos, escrever testes junto com o código e rodá-los a cada \emph{commit}.
      \end{dica}
      \begin{historia}
        Barry Boehm mediu esse efeito nos anos 1970 e 1980; os números variam entre projetos, mas a tendência se repete.
      \end{historia}
    \end{column}
  \end{columns}
\end{frame}

\begin{frame}{Verificação, validação e níveis de teste}
  \begin{columns}[c]
    \begin{column}{0.56\textwidth}
      \centering
      \begin{tikzpicture}[every node/.style={font=\scriptsize,minimum width=2.25cm,minimum height=6mm}]
        \node[caixa] (l1) at (1.15,0) {Requisitos};
        \node[caixa] (l2) at (1.55,-0.75) {Especificação};
        \node[caixa] (l3) at (1.95,-1.5) {Arquitetura};
        \node[caixa] (l4) at (2.35,-2.25) {Projeto detalhado};
        \node[caixaescura] (c) at (3.8,-3.0) {Código};
        \node[caixaambar] (r4) at (5.25,-2.25) {Teste de unidade};
        \node[caixaambar] (r3) at (5.65,-1.5) {Integração};
        \node[caixaambar] (r2) at (6.05,-0.75) {Teste de sistema};
        \node[caixaambar] (r1) at (6.45,0) {Aceitação};
        \foreach \a/\b in {l1/l2,l2/l3,l3/l4,l4/c,c/r4,r4/r3,r3/r2,r2/r1} \draw[seta] (\a) -- (\b);
        \foreach \a/\b in {l1/r1,l2/r2,l3/r3,l4/r4} \draw[pontilhado,<->] (\a) -- (\b);
      \end{tikzpicture}\\[2pt]
      {\tiny Modelo em V: cada nível de teste confere um nível de especificação.}
    \end{column}
    \begin{column}{0.42\textwidth}
      \begin{definicao}[Verificação]
        ``Estamos construindo o produto \textbf{corretamente}?'' Confere o software contra a especificação.
      \end{definicao}
      \begin{definicao}[Validação]
        ``Estamos construindo o produto \textbf{certo}?'' Confere se o software resolve o problema real do usuário.
      \end{definicao}
      \begin{dica}
        O teste de aceitação é feito com o cliente: no LabMonitor, os técnicos do laboratório usando o app.
      \end{dica}
    \end{column}
  \end{columns}
\end{frame}

\begin{frame}{Caixa-preta: partições e valores-limite}
  \begin{columns}[T]
    \begin{column}{0.5\textwidth}
      \begin{definicao}[Teste caixa-preta]
        Projetado a partir da \textbf{especificação}, sem olhar o código: entradas e saídas esperadas.
      \end{definicao}
      \begin{conceito}[Partição de equivalência]
        Divide as entradas em classes que o sistema deve tratar do mesmo jeito; basta um representante de cada classe.
      \end{conceito}
      \begin{conceito}[Análise de valor-limite]
        Defeitos se concentram nas bordas (\texttt{<} no lugar de \texttt{<=}). Teste o limite e seus vizinhos imediatos.
      \end{conceito}
    \end{column}
    \begin{column}{0.48\textwidth}
      \begin{exemplo}[Regra: a reserva dura de 1 a 4 horas inteiras]
        \centering
        \begin{tikzpicture}[x=0.62cm]
          \draw[cinzamedio,-{Stealth}] (-0.5,0) -- (6.5,0);
          \fill[vermelhoerro!25] (-0.5,-0.15) rectangle (0.5,0.15);
          \fill[verdeok!35] (0.5,-0.15) rectangle (4.5,0.15);
          \fill[vermelhoerro!25] (4.5,-0.15) rectangle (6.5,0.15);
          \foreach \v in {0,1,2,3,4,5,6} \node[font=\tiny,below] at (\v,-0.15) {\v};
          \node[font=\tiny,above] at (0,0.15) {inválida};
          \node[font=\tiny,above] at (2.5,0.15) {válida};
          \node[font=\tiny,above] at (5.5,0.15) {inválida};
        \end{tikzpicture}\par
        \scriptsize
        \begin{tabular}{@{}l l@{}}
          \toprule
          \cabecalho{Casos} & \cabecalho{Por quê} \\
          \midrule
          2 h & representante da classe válida \\
          0 h e 1 h & limite inferior e vizinho \\
          4 h e 5 h & limite superior e vizinho \\
          $-1$ h, ``abc'' & entradas absurdas \\
          \bottomrule
        \end{tabular}
      \end{exemplo}
    \end{column}
  \end{columns}
\end{frame}

\begin{frame}[fragile]{Caixa-branca: critérios de cobertura}
  \begin{columns}[T]
    \begin{column}{0.46\textwidth}
      \begin{definicao}[Teste caixa-branca]
        Projetado a partir da \textbf{estrutura interna} do código: comandos, decisões e caminhos.
      \end{definicao}
\begin{codigo}[title=Pseudocódigo]{bash}
se temp >= limite E alertaAtivo:
    notificar()
registrar(temp)
\end{codigo}
      \begin{atencao}
        Cobertura mede o que foi \textbf{executado}, não o que foi \textbf{conferido}: um teste sem \cod{expect} também ``cobre''.
      \end{atencao}
    \end{column}
    \begin{column}{0.52\textwidth}
      \scriptsize
      \begin{tabularx}{\linewidth}{@{}p{2.1cm} L@{}}
        \toprule
        \cabecalho{Critério} & \cabecalho{Exige, no exemplo} \\
        \midrule
        \textbf{Comandos} & cada linha rodar ao menos uma vez: 1 teste (30\,°C, ativo) basta \\
        \textbf{Decisões} & cada \cod{se} (\emph{branch}) ser verdadeiro e falso: 2 testes \\
        \textbf{Condições} & cada condição simples verdadeira e falsa: \cod{temp} e \cod{alertaAtivo} variando \\
        \textbf{Caminhos} & todas as rotas possíveis; cresce exponencialmente com laços e \cod{se}s \\
        \bottomrule
      \end{tabularx}
      \begin{conceito}[Hierarquia]
        Caminhos $\Rightarrow$ decisões $\Rightarrow$ comandos: quem satisfaz um critério mais forte satisfaz os mais fracos.
      \end{conceito}
    \end{column}
  \end{columns}
\end{frame}

\begin{frame}{Dublês de teste}
  \begin{columns}[T]
    \begin{column}{0.62\textwidth}
      \scriptsize
      \begin{tabularx}{\linewidth}{@{}l L L@{}}
        \toprule
        \cabecalho{Dublê} & \cabecalho{O que é} & \cabecalho{Exemplo no LabMonitor} \\
        \midrule
        \textbf{Dummy} & objeto só para preencher um parâmetro; nunca é usado & um \emph{logger} vazio exigido pelo construtor \\
        \textbf{Stub} & devolve respostas prontas & repositório que sempre devolve 3 bancadas \\
        \textbf{Spy} & stub que também registra as chamadas & conta quantas notificações foram enviadas \\
        \textbf{Mock} & objeto com \textbf{expectativas} de interação, verificadas no teste & ``\cod{reservar(7)} deve ser chamado uma vez'' \\
        \textbf{Fake} & implementação funcional simplificada & banco em memória, API falsa com dados fixos \\
        \bottomrule
      \end{tabularx}
    \end{column}
    \begin{column}{0.36\textwidth}
      \begin{definicao}[Dublê de teste]
        Objeto que substitui uma dependência real (rede, banco, sensor, relógio) durante o teste.
      \end{definicao}
      \begin{analogia}
        Como o dublê de cinema: entra nas cenas perigosas ou caras no lugar do ator.
      \end{analogia}
      \begin{historia}
        A classificação é de Gerard Meszaros, no livro \emph{xUnit Test Patterns} (2007).
      \end{historia}
    \end{column}
  \end{columns}
\end{frame}

\begin{frame}{Propriedades de um bom teste}
  \begin{columns}[T]
    \begin{column}{0.5\textwidth}
      \scriptsize
      \begin{tabularx}{\linewidth}{@{}l L@{}}
        \toprule
        \cabecalho{FIRST} & \cabecalho{Significado} \\
        \midrule
        \textbf{F}ast & rápido: roda em milissegundos, então roda sempre \\
        \textbf{I}ndependent & não depende da ordem nem de outros testes \\
        \textbf{R}epeatable & mesmo resultado em qualquer máquina e horário \\
        \textbf{S}elf-validating & passa ou falha sozinho, sem alguém olhar a saída \\
        \textbf{T}imely & escrito junto com o código (ou antes) \\
        \bottomrule
      \end{tabularx}
      \begin{conceito}[Estrutura]
        \textbf{AAA}: \emph{Arrange}, \emph{Act}, \emph{Assert}.\\
        \textbf{BDD}: \emph{Given}, \emph{When}, \emph{Then} (dado, quando, então).
      \end{conceito}
    \end{column}
    \begin{column}{0.48\textwidth}
      \begin{erro}[Teste instável (\emph{flaky})]
        Passa e falha sem o código mudar. Causas típicas: depender do relógio real, da rede, da ordem de execução, de esperas fixas (\emph{sleep}) ou de estado compartilhado entre testes.
      \end{erro}
      \begin{dica}[Teste também é código]
        Merece nomes claros, nada de duplicação e revisão. Um bom nome de teste descreve uma \textbf{regra}: ``reserva sobreposta é recusada''.
      \end{dica}
    \end{column}
  \end{columns}
\end{frame}

\begin{frame}[fragile]{TDD, BDD e regressão}
  \begin{columns}[T]
    \begin{column}{0.4\textwidth}
      \centering
      \begin{tikzpicture}
        \node[circle,fill=vermelhoerro,text=white,font=\scriptsize\bfseries,minimum size=1.35cm,align=center] (v) at (90:1.35) {Vermelho\\ \tiny teste falha};
        \node[circle,fill=verdeok,text=white,font=\scriptsize\bfseries,minimum size=1.35cm,align=center] (g) at (-30:1.35) {Verde\\ \tiny passa};
        \node[circle,fill=azulpetroleo,text=white,font=\scriptsize\bfseries,minimum size=1.35cm,align=center] (r) at (210:1.35) {Refatorar\\ \tiny melhora};
        \draw[setaambar] (v) to[bend left=25] (g);
        \draw[setaambar] (g) to[bend left=25] (r);
        \draw[setaambar] (r) to[bend left=25] (v);
      \end{tikzpicture}
      \begin{definicao}[TDD]
        \emph{Test-Driven Development} (Kent Beck): o teste é escrito \textbf{antes} do código e guia o projeto.
      \end{definicao}
    \end{column}
    \begin{column}{0.58\textwidth}
\begin{codigo}[title=BDD: cenário em linguagem natural (Gherkin)]{bash}
Funcionalidade: Reserva de bancada
  Cenário: horário já ocupado
    Dado que a bancada B7 está reservada das 14h às 16h
    Quando eu peço a B7 das 15h às 17h
    Então vejo "Horário indisponível"
\end{codigo}
      \begin{definicao}[BDD]
        \emph{Behavior-Driven Development}: descreve o comportamento em exemplos que clientes, testadores e devs entendem; os cenários viram testes.
      \end{definicao}
      \begin{definicao}[Teste de regressão]
        Rodar de novo os testes existentes após cada mudança, para garantir que nada que funcionava quebrou. Automatizar é o que torna isso viável.
      \end{definicao}
    \end{column}
  \end{columns}
\end{frame}

\begin{frame}{Desafios do teste em apps móveis}
  \begin{columns}[T]
    \begin{column}{0.5\textwidth}
      \scriptsize
      \begin{tabularx}{\linewidth}{@{}p{2.0cm} L@{}}
        \toprule
        \cabecalho{Desafio} & \cabecalho{O que pode dar errado} \\
        \midrule
        Fragmentação & milhares de modelos Android, versões de sistema, tamanhos e densidades de tela \\
        Rede & sem sinal, 3G lento, troca Wi-Fi/móvel no meio de uma requisição \\
        Interrupções & ligação, notificação, app enviado ao fundo e encerrado pelo sistema \\
        Hardware & câmera, GPS, Bluetooth e sensores diferentes ou ausentes \\
        Permissões & usuário nega, concede ``só desta vez'' ou revoga depois \\
        Recursos & bateria fraca, pouca memória, armazenamento cheio \\
        Pessoas & idioma, fonte grande, leitor de tela, modo escuro \\
        \bottomrule
      \end{tabularx}
    \end{column}
    \begin{column}{0.48\textwidth}
      \begin{conceito}[Estratégia]
        \begin{itemize}
          \item Automatizar a lógica em testes de unidade e widget (rápidos, sem aparelho).
          \item Definir uma \textbf{matriz de dispositivos} representativa (versões antigas e novas, tela pequena e grande).
          \item Usar emuladores no dia a dia e aparelhos reais ou \emph{device farms} antes de publicar.
          \item Simular rede ruim e interrupções.
        \end{itemize}
      \end{conceito}
      \begin{dica}
        As estatísticas do Android Studio e do console da loja mostram quais versões seus usuários realmente usam.
      \end{dica}
    \end{column}
  \end{columns}
\end{frame}

\begin{frame}{Qualidade é mais do que testar}
  \begin{columns}[T]
    \begin{column}{0.5\textwidth}
      \begin{definicao}[Modelo de qualidade ISO/IEC 25010]
        Organiza a qualidade de um produto de software em características como adequação funcional, eficiência de desempenho, compatibilidade, usabilidade, confiabilidade, segurança, manutenibilidade e portabilidade.
      \end{definicao}
      \begin{dica}
        A revisão de 2023 renomeou algumas características (usabilidade virou capacidade de interação) e incluiu \emph{safety}.
      \end{dica}
    \end{column}
    \begin{column}{0.48\textwidth}
      \begin{conceito}[Técnicas que se somam]
        \begin{itemize}
          \item \textbf{Estáticas} (sem executar): revisão de código, análise estática, \emph{lint}, tipagem forte.
          \item \textbf{Dinâmicas} (executando): testes automatizados e exploratórios.
          \item \textbf{Em produção}: monitoramento de falhas, métricas e \emph{feedback} dos usuários.
        \end{itemize}
      \end{conceito}
      \begin{destaque}
        Qualidade se \textbf{constrói} durante o desenvolvimento; testes confirmam que ela está lá.
      \end{destaque}
    \end{column}
  \end{columns}
\end{frame}

% =============================================================================
\section{Fundamentos de teste de software}

\begin{frame}{Erro, defeito e falha}
  \begin{definicao}[Três palavras que não são sinônimos]
    \textbf{Erro} (\emph{error}): engano cometido por uma pessoa.
    \textbf{Defeito} (\emph{defect}, \emph{bug}): imperfeição no código ou em outro artefato, consequência do erro.
    \textbf{Falha} (\emph{failure}): comportamento incorreto observado quando o defeito é executado.
  \end{definicao}
  \centering
  \begin{tikzpicture}[node distance=17mm]
    \node[caixa,minimum width=3.0cm] (e) {\textbf{Erro}\\ \tiny o dev usa \texttt{>} em vez de \texttt{>=}};
    \node[caixapetroleo,minimum width=3.0cm,right=of e] (d) {\textbf{Defeito}\\ \tiny a regra de alerta fica errada};
    \node[caixaambar,minimum width=3.2cm,right=of d] (f) {\textbf{Falha}\\ \tiny a 30,0\,°C nenhum alerta aparece};
    \draw[seta] (e) -- node[rotulo,above]{\tiny introduz} (d);
    \draw[seta] (d) -- node[rotulo,above]{\tiny executado} (f);
  \end{tikzpicture}
  \vspace{2pt}
  \begin{columns}[T]
    \begin{column}{0.48\textwidth}
      \begin{conceito}[Consequência]
        Um defeito pode ficar anos sem causar falha, se a linha nunca for executada nas condições certas. Testar é \textbf{provocar} essas condições.
      \end{conceito}
    \end{column}
    \begin{column}{0.48\textwidth}
      \begin{historia}
        Em 1947, a equipe do computador Harvard Mark II, da qual Grace Hopper fazia parte, achou uma mariposa presa num relé e a colou no diário de bordo. A palavra \emph{bug} já existia, mas o episódio a popularizou.
      \end{historia}
    \end{column}
  \end{columns}
\end{frame}

\begin{frame}{Para que serve testar}
  \begin{columns}[T]
    \begin{column}{0.5\textwidth}
      \begin{conceito}[Objetivos dos testes]
        \begin{itemize}
          \item Encontrar defeitos antes do usuário.
          \item Dar confiança para entregar.
          \item Impedir que defeitos corrigidos voltem (\textbf{regressão}).
          \item Documentar, em código executável, como o sistema deve se comportar.
          \item Permitir \textbf{refatorar} sem medo.
        \end{itemize}
      \end{conceito}
      \begin{atencao}[Princípio fundamental]
        Testes mostram que defeitos \textbf{existem}; nunca provam que \textbf{não existem}. A ideia é de Edsger Dijkstra.
      \end{atencao}
    \end{column}
    \begin{column}{0.48\textwidth}
      \begin{pordentro}[Teste exaustivo é impossível]
        Uma função com dois inteiros de 32 bits tem $2^{64}$ combinações de entrada, cerca de $1{,}8\times10^{19}$. A um bilhão de testes por segundo, seriam quase 600 anos. Por isso escolhemos casos \textbf{representativos} com técnicas de projeto de testes.
      \end{pordentro}
      \begin{analogia}
        O \emph{crash test} de um carro não bate o veículo em todas as velocidades e ângulos: escolhe poucos cenários que representam os acidentes reais.
      \end{analogia}
    \end{column}
  \end{columns}
\end{frame}

\begin{frame}{O custo de achar tarde}
  \begin{columns}[c]
    \begin{column}{0.52\textwidth}
      \centering
      \begin{tikzpicture}[x=0.95cm,y=0.5cm]
        \draw[cinzamedio] (0,0) -- (5.3,0);
        \foreach \x/\h/\t in {0.4/0.6/Requisitos,1.4/1.2/Projeto,2.4/2.2/Código,3.4/3.8/Teste,4.4/6.4/Produção}{
          \fill[azulpetroleo] (\x,0) rectangle ++(0.6,\h);
          \node[font=\tiny,rotate=35,anchor=north east] at (\x+0.45,-0.1) {\t};}
        \fill[ambar] (4.4,0) rectangle ++(0.6,6.4);
        \node[rotulo,font=\tiny,align=center] at (2.2,6.5) {custo relativo de corrigir\\ (ilustrativo)};
      \end{tikzpicture}
    \end{column}
    \begin{column}{0.46\textwidth}
      \begin{conceito}[Por que custa mais]
        Um defeito achado em produção exige investigar relatos, reproduzir, corrigir, testar de novo, gerar versão, passar pela revisão da loja e esperar os usuários atualizarem. No app, ainda há dados já gravados errado.
      \end{conceito}
      \begin{dica}[\emph{Shift left}]
        Antecipar a qualidade: revisar requisitos, escrever testes junto com o código e rodá-los a cada \emph{commit}.
      \end{dica}
      \begin{historia}
        Barry Boehm mediu esse efeito nos anos 1970 e 1980; os números variam entre projetos, mas a tendência se repete.
      \end{historia}
    \end{column}
  \end{columns}
\end{frame}

\begin{frame}{Verificação, validação e níveis de teste}
  \begin{columns}[c]
    \begin{column}{0.56\textwidth}
      \centering
      \begin{tikzpicture}[every node/.style={font=\scriptsize,minimum width=2.25cm,minimum height=6mm}]
        \node[caixa] (l1) at (1.15,0) {Requisitos};
        \node[caixa] (l2) at (1.55,-0.75) {Especificação};
        \node[caixa] (l3) at (1.95,-1.5) {Arquitetura};
        \node[caixa] (l4) at (2.35,-2.25) {Projeto detalhado};
        \node[caixaescura] (c) at (3.8,-3.0) {Código};
        \node[caixaambar] (r4) at (5.25,-2.25) {Teste de unidade};
        \node[caixaambar] (r3) at (5.65,-1.5) {Integração};
        \node[caixaambar] (r2) at (6.05,-0.75) {Teste de sistema};
        \node[caixaambar] (r1) at (6.45,0) {Aceitação};
        \foreach \a/\b in {l1/l2,l2/l3,l3/l4,l4/c,c/r4,r4/r3,r3/r2,r2/r1} \draw[seta] (\a) -- (\b);
        \foreach \a/\b in {l1/r1,l2/r2,l3/r3,l4/r4} \draw[pontilhado,<->] (\a) -- (\b);
      \end{tikzpicture}\\[2pt]
      {\tiny Modelo em V: cada nível de teste confere um nível de especificação.}
    \end{column}
    \begin{column}{0.42\textwidth}
      \begin{definicao}[Verificação]
        ``Estamos construindo o produto \textbf{corretamente}?'' Confere o software contra a especificação.
      \end{definicao}
      \begin{definicao}[Validação]
        ``Estamos construindo o produto \textbf{certo}?'' Confere se o software resolve o problema real do usuário.
      \end{definicao}
      \begin{dica}
        O teste de aceitação é feito com o cliente: no LabMonitor, os técnicos do laboratório usando o app.
      \end{dica}
    \end{column}
  \end{columns}
\end{frame}

\begin{frame}{Caixa-preta: partições e valores-limite}
  \begin{columns}[T]
    \begin{column}{0.5\textwidth}
      \begin{definicao}[Teste caixa-preta]
        Projetado a partir da \textbf{especificação}, sem olhar o código: entradas e saídas esperadas.
      \end{definicao}
      \begin{conceito}[Partição de equivalência]
        Divide as entradas em classes que o sistema deve tratar do mesmo jeito; basta um representante de cada classe.
      \end{conceito}
      \begin{conceito}[Análise de valor-limite]
        Defeitos se concentram nas bordas (\texttt{<} no lugar de \texttt{<=}). Teste o limite e seus vizinhos imediatos.
      \end{conceito}
    \end{column}
    \begin{column}{0.48\textwidth}
      \begin{exemplo}[Regra: a reserva dura de 1 a 4 horas inteiras]
        \centering
        \begin{tikzpicture}[x=0.62cm]
          \draw[cinzamedio,-{Stealth}] (-0.5,0) -- (6.5,0);
          \fill[vermelhoerro!25] (-0.5,-0.15) rectangle (0.5,0.15);
          \fill[verdeok!35] (0.5,-0.15) rectangle (4.5,0.15);
          \fill[vermelhoerro!25] (4.5,-0.15) rectangle (6.5,0.15);
          \foreach \v in {0,1,2,3,4,5,6} \node[font=\tiny,below] at (\v,-0.15) {\v};
          \node[font=\tiny,above] at (0,0.15) {inválida};
          \node[font=\tiny,above] at (2.5,0.15) {válida};
          \node[font=\tiny,above] at (5.5,0.15) {inválida};
        \end{tikzpicture}\par
        \scriptsize
        \begin{tabular}{@{}l l@{}}
          \toprule
          \cabecalho{Casos} & \cabecalho{Por quê} \\
          \midrule
          2 h & representante da classe válida \\
          0 h e 1 h & limite inferior e vizinho \\
          4 h e 5 h & limite superior e vizinho \\
          $-1$ h, ``abc'' & entradas absurdas \\
          \bottomrule
        \end{tabular}
      \end{exemplo}
    \end{column}
  \end{columns}
\end{frame}

\begin{frame}[fragile]{Caixa-branca: critérios de cobertura}
  \begin{columns}[T]
    \begin{column}{0.46\textwidth}
      \begin{definicao}[Teste caixa-branca]
        Projetado a partir da \textbf{estrutura interna} do código: comandos, decisões e caminhos.
      \end{definicao}
\begin{codigo}[title=Pseudocódigo]{bash}
se temp >= limite E alertaAtivo:
    notificar()
registrar(temp)
\end{codigo}
      \begin{atencao}
        Cobertura mede o que foi \textbf{executado}, não o que foi \textbf{conferido}: um teste sem \cod{expect} também ``cobre''.
      \end{atencao}
    \end{column}
    \begin{column}{0.52\textwidth}
      \scriptsize
      \begin{tabularx}{\linewidth}{@{}p{2.1cm} L@{}}
        \toprule
        \cabecalho{Critério} & \cabecalho{Exige, no exemplo} \\
        \midrule
        \textbf{Comandos} & cada linha rodar ao menos uma vez: 1 teste (30\,°C, ativo) basta \\
        \textbf{Decisões} & cada \cod{se} (\emph{branch}) ser verdadeiro e falso: 2 testes \\
        \textbf{Condições} & cada condição simples verdadeira e falsa: \cod{temp} e \cod{alertaAtivo} variando \\
        \textbf{Caminhos} & todas as rotas possíveis; cresce exponencialmente com laços e \cod{se}s \\
        \bottomrule
      \end{tabularx}
      \begin{conceito}[Hierarquia]
        Caminhos $\Rightarrow$ decisões $\Rightarrow$ comandos: quem satisfaz um critério mais forte satisfaz os mais fracos.
      \end{conceito}
    \end{column}
  \end{columns}
\end{frame}

\begin{frame}{Dublês de teste}
  \begin{columns}[T]
    \begin{column}{0.62\textwidth}
      \scriptsize
      \begin{tabularx}{\linewidth}{@{}l L L@{}}
        \toprule
        \cabecalho{Dublê} & \cabecalho{O que é} & \cabecalho{Exemplo no LabMonitor} \\
        \midrule
        \textbf{Dummy} & objeto só para preencher um parâmetro; nunca é usado & um \emph{logger} vazio exigido pelo construtor \\
        \textbf{Stub} & devolve respostas prontas & repositório que sempre devolve 3 bancadas \\
        \textbf{Spy} & stub que também registra as chamadas & conta quantas notificações foram enviadas \\
        \textbf{Mock} & objeto com \textbf{expectativas} de interação, verificadas no teste & ``\cod{reservar(7)} deve ser chamado uma vez'' \\
        \textbf{Fake} & implementação funcional simplificada & banco em memória, API falsa com dados fixos \\
        \bottomrule
      \end{tabularx}
    \end{column}
    \begin{column}{0.36\textwidth}
      \begin{definicao}[Dublê de teste]
        Objeto que substitui uma dependência real (rede, banco, sensor, relógio) durante o teste.
      \end{definicao}
      \begin{analogia}
        Como o dublê de cinema: entra nas cenas perigosas ou caras no lugar do ator.
      \end{analogia}
      \begin{historia}
        A classificação é de Gerard Meszaros, no livro \emph{xUnit Test Patterns} (2007).
      \end{historia}
    \end{column}
  \end{columns}
\end{frame}

\begin{frame}{Propriedades de um bom teste}
  \begin{columns}[T]
    \begin{column}{0.5\textwidth}
      \scriptsize
      \begin{tabularx}{\linewidth}{@{}l L@{}}
        \toprule
        \cabecalho{FIRST} & \cabecalho{Significado} \\
        \midrule
        \textbf{F}ast & rápido: roda em milissegundos, então roda sempre \\
        \textbf{I}ndependent & não depende da ordem nem de outros testes \\
        \textbf{R}epeatable & mesmo resultado em qualquer máquina e horário \\
        \textbf{S}elf-validating & passa ou falha sozinho, sem alguém olhar a saída \\
        \textbf{T}imely & escrito junto com o código (ou antes) \\
        \bottomrule
      \end{tabularx}
      \begin{conceito}[Estrutura]
        \textbf{AAA}: \emph{Arrange}, \emph{Act}, \emph{Assert}.\\
        \textbf{BDD}: \emph{Given}, \emph{When}, \emph{Then} (dado, quando, então).
      \end{conceito}
    \end{column}
    \begin{column}{0.48\textwidth}
      \begin{erro}[Teste instável (\emph{flaky})]
        Passa e falha sem o código mudar. Causas típicas: depender do relógio real, da rede, da ordem de execução, de esperas fixas (\emph{sleep}) ou de estado compartilhado entre testes.
      \end{erro}
      \begin{dica}[Teste também é código]
        Merece nomes claros, nada de duplicação e revisão. Um bom nome de teste descreve uma \textbf{regra}: ``reserva sobreposta é recusada''.
      \end{dica}
    \end{column}
  \end{columns}
\end{frame}

\begin{frame}[fragile]{TDD, BDD e regressão}
  \begin{columns}[T]
    \begin{column}{0.4\textwidth}
      \centering
      \begin{tikzpicture}
        \node[circle,fill=vermelhoerro,text=white,font=\scriptsize\bfseries,minimum size=1.35cm,align=center] (v) at (90:1.35) {Vermelho\\ \tiny teste falha};
        \node[circle,fill=verdeok,text=white,font=\scriptsize\bfseries,minimum size=1.35cm,align=center] (g) at (-30:1.35) {Verde\\ \tiny passa};
        \node[circle,fill=azulpetroleo,text=white,font=\scriptsize\bfseries,minimum size=1.35cm,align=center] (r) at (210:1.35) {Refatorar\\ \tiny melhora};
        \draw[setaambar] (v) to[bend left=25] (g);
        \draw[setaambar] (g) to[bend left=25] (r);
        \draw[setaambar] (r) to[bend left=25] (v);
      \end{tikzpicture}
      \begin{definicao}[TDD]
        \emph{Test-Driven Development} (Kent Beck): o teste é escrito \textbf{antes} do código e guia o projeto.
      \end{definicao}
    \end{column}
    \begin{column}{0.58\textwidth}
\begin{codigo}[title=BDD: cenário em linguagem natural (Gherkin)]{bash}
Funcionalidade: Reserva de bancada
  Cenário: horário já ocupado
    Dado que a bancada B7 está reservada das 14h às 16h
    Quando eu peço a B7 das 15h às 17h
    Então vejo "Horário indisponível"
\end{codigo}
      \begin{definicao}[BDD]
        \emph{Behavior-Driven Development}: descreve o comportamento em exemplos que clientes, testadores e devs entendem; os cenários viram testes.
      \end{definicao}
      \begin{definicao}[Teste de regressão]
        Rodar de novo os testes existentes após cada mudança, para garantir que nada que funcionava quebrou. Automatizar é o que torna isso viável.
      \end{definicao}
    \end{column}
  \end{columns}
\end{frame}

\begin{frame}{Desafios do teste em apps móveis}
  \begin{columns}[T]
    \begin{column}{0.5\textwidth}
      \scriptsize
      \begin{tabularx}{\linewidth}{@{}p{2.0cm} L@{}}
        \toprule
        \cabecalho{Desafio} & \cabecalho{O que pode dar errado} \\
        \midrule
        Fragmentação & milhares de modelos Android, versões de sistema, tamanhos e densidades de tela \\
        Rede & sem sinal, 3G lento, troca Wi-Fi/móvel no meio de uma requisição \\
        Interrupções & ligação, notificação, app enviado ao fundo e encerrado pelo sistema \\
        Hardware & câmera, GPS, Bluetooth e sensores diferentes ou ausentes \\
        Permissões & usuário nega, concede ``só desta vez'' ou revoga depois \\
        Recursos & bateria fraca, pouca memória, armazenamento cheio \\
        Pessoas & idioma, fonte grande, leitor de tela, modo escuro \\
        \bottomrule
      \end{tabularx}
    \end{column}
    \begin{column}{0.48\textwidth}
      \begin{conceito}[Estratégia]
        \begin{itemize}
          \item Automatizar a lógica em testes de unidade e widget (rápidos, sem aparelho).
          \item Definir uma \textbf{matriz de dispositivos} representativa (versões antigas e novas, tela pequena e grande).
          \item Usar emuladores no dia a dia e aparelhos reais ou \emph{device farms} antes de publicar.
          \item Simular rede ruim e interrupções.
        \end{itemize}
      \end{conceito}
      \begin{dica}
        As estatísticas do Android Studio e do console da loja mostram quais versões seus usuários realmente usam.
      \end{dica}
    \end{column}
  \end{columns}
\end{frame}

\begin{frame}{Qualidade é mais do que testar}
  \begin{columns}[T]
    \begin{column}{0.5\textwidth}
      \begin{definicao}[Modelo de qualidade ISO/IEC 25010]
        Organiza a qualidade de um produto de software em características como adequação funcional, eficiência de desempenho, compatibilidade, usabilidade, confiabilidade, segurança, manutenibilidade e portabilidade.
      \end{definicao}
      \begin{dica}
        A revisão de 2023 renomeou algumas características (usabilidade virou capacidade de interação) e incluiu \emph{safety}.
      \end{dica}
    \end{column}
    \begin{column}{0.48\textwidth}
      \begin{conceito}[Técnicas que se somam]
        \begin{itemize}
          \item \textbf{Estáticas} (sem executar): revisão de código, análise estática, \emph{lint}, tipagem forte.
          \item \textbf{Dinâmicas} (executando): testes automatizados e exploratórios.
          \item \textbf{Em produção}: monitoramento de falhas, métricas e \emph{feedback} dos usuários.
        \end{itemize}
      \end{conceito}
      \begin{destaque}
        Qualidade se \textbf{constrói} durante o desenvolvimento; testes confirmam que ela está lá.
      \end{destaque}
    \end{column}
  \end{columns}
\end{frame}

% =============================================================================
\section{Fundamentos de teste de software}

\begin{frame}{Erro, defeito e falha}
  \begin{definicao}[Três palavras que não são sinônimos]
    \textbf{Erro} (\emph{error}): engano cometido por uma pessoa.
    \textbf{Defeito} (\emph{defect}, \emph{bug}): imperfeição no código ou em outro artefato, consequência do erro.
    \textbf{Falha} (\emph{failure}): comportamento incorreto observado quando o defeito é executado.
  \end{definicao}
  \centering
  \begin{tikzpicture}[node distance=17mm]
    \node[caixa,minimum width=3.0cm] (e) {\textbf{Erro}\\ \tiny o dev usa \texttt{>} em vez de \texttt{>=}};
    \node[caixapetroleo,minimum width=3.0cm,right=of e] (d) {\textbf{Defeito}\\ \tiny a regra de alerta fica errada};
    \node[caixaambar,minimum width=3.2cm,right=of d] (f) {\textbf{Falha}\\ \tiny a 30,0\,°C nenhum alerta aparece};
    \draw[seta] (e) -- node[rotulo,above]{\tiny introduz} (d);
    \draw[seta] (d) -- node[rotulo,above]{\tiny executado} (f);
  \end{tikzpicture}
  \vspace{2pt}
  \begin{columns}[T]
    \begin{column}{0.48\textwidth}
      \begin{conceito}[Consequência]
        Um defeito pode ficar anos sem causar falha, se a linha nunca for executada nas condições certas. Testar é \textbf{provocar} essas condições.
      \end{conceito}
    \end{column}
    \begin{column}{0.48\textwidth}
      \begin{historia}
        Em 1947, a equipe do computador Harvard Mark II, da qual Grace Hopper fazia parte, achou uma mariposa presa num relé e a colou no diário de bordo. A palavra \emph{bug} já existia, mas o episódio a popularizou.
      \end{historia}
    \end{column}
  \end{columns}
\end{frame}

\begin{frame}{Para que serve testar}
  \begin{columns}[T]
    \begin{column}{0.5\textwidth}
      \begin{conceito}[Objetivos dos testes]
        \begin{itemize}
          \item Encontrar defeitos antes do usuário.
          \item Dar confiança para entregar.
          \item Impedir que defeitos corrigidos voltem (\textbf{regressão}).
          \item Documentar, em código executável, como o sistema deve se comportar.
          \item Permitir \textbf{refatorar} sem medo.
        \end{itemize}
      \end{conceito}
      \begin{atencao}[Princípio fundamental]
        Testes mostram que defeitos \textbf{existem}; nunca provam que \textbf{não existem}. A ideia é de Edsger Dijkstra.
      \end{atencao}
    \end{column}
    \begin{column}{0.48\textwidth}
      \begin{pordentro}[Teste exaustivo é impossível]
        Uma função com dois inteiros de 32 bits tem $2^{64}$ combinações de entrada, cerca de $1{,}8\times10^{19}$. A um bilhão de testes por segundo, seriam quase 600 anos. Por isso escolhemos casos \textbf{representativos} com técnicas de projeto de testes.
      \end{pordentro}
      \begin{analogia}
        O \emph{crash test} de um carro não bate o veículo em todas as velocidades e ângulos: escolhe poucos cenários que representam os acidentes reais.
      \end{analogia}
    \end{column}
  \end{columns}
\end{frame}

\begin{frame}{O custo de achar tarde}
  \begin{columns}[c]
    \begin{column}{0.52\textwidth}
      \centering
      \begin{tikzpicture}[x=0.95cm,y=0.5cm]
        \draw[cinzamedio] (0,0) -- (5.3,0);
        \foreach \x/\h/\t in {0.4/0.6/Requisitos,1.4/1.2/Projeto,2.4/2.2/Código,3.4/3.8/Teste,4.4/6.4/Produção}{
          \fill[azulpetroleo] (\x,0) rectangle ++(0.6,\h);
          \node[font=\tiny,rotate=35,anchor=north east] at (\x+0.45,-0.1) {\t};}
        \fill[ambar] (4.4,0) rectangle ++(0.6,6.4);
        \node[rotulo,font=\tiny,align=center] at (2.2,6.5) {custo relativo de corrigir\\ (ilustrativo)};
      \end{tikzpicture}
    \end{column}
    \begin{column}{0.46\textwidth}
      \begin{conceito}[Por que custa mais]
        Um defeito achado em produção exige investigar relatos, reproduzir, corrigir, testar de novo, gerar versão, passar pela revisão da loja e esperar os usuários atualizarem. No app, ainda há dados já gravados errado.
      \end{conceito}
      \begin{dica}[\emph{Shift left}]
        Antecipar a qualidade: revisar requisitos, escrever testes junto com o código e rodá-los a cada \emph{commit}.
      \end{dica}
      \begin{historia}
        Barry Boehm mediu esse efeito nos anos 1970 e 1980; os números variam entre projetos, mas a tendência se repete.
      \end{historia}
    \end{column}
  \end{columns}
\end{frame}

\begin{frame}{Verificação, validação e níveis de teste}
  \begin{columns}[c]
    \begin{column}{0.56\textwidth}
      \centering
      \begin{tikzpicture}[every node/.style={font=\scriptsize,minimum width=2.25cm,minimum height=6mm}]
        \node[caixa] (l1) at (1.15,0) {Requisitos};
        \node[caixa] (l2) at (1.55,-0.75) {Especificação};
        \node[caixa] (l3) at (1.95,-1.5) {Arquitetura};
        \node[caixa] (l4) at (2.35,-2.25) {Projeto detalhado};
        \node[caixaescura] (c) at (3.8,-3.0) {Código};
        \node[caixaambar] (r4) at (5.25,-2.25) {Teste de unidade};
        \node[caixaambar] (r3) at (5.65,-1.5) {Integração};
        \node[caixaambar] (r2) at (6.05,-0.75) {Teste de sistema};
        \node[caixaambar] (r1) at (6.45,0) {Aceitação};
        \foreach \a/\b in {l1/l2,l2/l3,l3/l4,l4/c,c/r4,r4/r3,r3/r2,r2/r1} \draw[seta] (\a) -- (\b);
        \foreach \a/\b in {l1/r1,l2/r2,l3/r3,l4/r4} \draw[pontilhado,<->] (\a) -- (\b);
      \end{tikzpicture}\\[2pt]
      {\tiny Modelo em V: cada nível de teste confere um nível de especificação.}
    \end{column}
    \begin{column}{0.42\textwidth}
      \begin{definicao}[Verificação]
        ``Estamos construindo o produto \textbf{corretamente}?'' Confere o software contra a especificação.
      \end{definicao}
      \begin{definicao}[Validação]
        ``Estamos construindo o produto \textbf{certo}?'' Confere se o software resolve o problema real do usuário.
      \end{definicao}
      \begin{dica}
        O teste de aceitação é feito com o cliente: no LabMonitor, os técnicos do laboratório usando o app.
      \end{dica}
    \end{column}
  \end{columns}
\end{frame}

\begin{frame}{Caixa-preta: partições e valores-limite}
  \begin{columns}[T]
    \begin{column}{0.5\textwidth}
      \begin{definicao}[Teste caixa-preta]
        Projetado a partir da \textbf{especificação}, sem olhar o código: entradas e saídas esperadas.
      \end{definicao}
      \begin{conceito}[Partição de equivalência]
        Divide as entradas em classes que o sistema deve tratar do mesmo jeito; basta um representante de cada classe.
      \end{conceito}
      \begin{conceito}[Análise de valor-limite]
        Defeitos se concentram nas bordas (\texttt{<} no lugar de \texttt{<=}). Teste o limite e seus vizinhos imediatos.
      \end{conceito}
    \end{column}
    \begin{column}{0.48\textwidth}
      \begin{exemplo}[Regra: a reserva dura de 1 a 4 horas inteiras]
        \centering
        \begin{tikzpicture}[x=0.62cm]
          \draw[cinzamedio,-{Stealth}] (-0.5,0) -- (6.5,0);
          \fill[vermelhoerro!25] (-0.5,-0.15) rectangle (0.5,0.15);
          \fill[verdeok!35] (0.5,-0.15) rectangle (4.5,0.15);
          \fill[vermelhoerro!25] (4.5,-0.15) rectangle (6.5,0.15);
          \foreach \v in {0,1,2,3,4,5,6} \node[font=\tiny,below] at (\v,-0.15) {\v};
          \node[font=\tiny,above] at (0,0.15) {inválida};
          \node[font=\tiny,above] at (2.5,0.15) {válida};
          \node[font=\tiny,above] at (5.5,0.15) {inválida};
        \end{tikzpicture}\par
        \scriptsize
        \begin{tabular}{@{}l l@{}}
          \toprule
          \cabecalho{Casos} & \cabecalho{Por quê} \\
          \midrule
          2 h & representante da classe válida \\
          0 h e 1 h & limite inferior e vizinho \\
          4 h e 5 h & limite superior e vizinho \\
          $-1$ h, ``abc'' & entradas absurdas \\
          \bottomrule
        \end{tabular}
      \end{exemplo}
    \end{column}
  \end{columns}
\end{frame}

\begin{frame}[fragile]{Caixa-branca: critérios de cobertura}
  \begin{columns}[T]
    \begin{column}{0.46\textwidth}
      \begin{definicao}[Teste caixa-branca]
        Projetado a partir da \textbf{estrutura interna} do código: comandos, decisões e caminhos.
      \end{definicao}
\begin{codigo}[title=Pseudocódigo]{bash}
se temp >= limite E alertaAtivo:
    notificar()
registrar(temp)
\end{codigo}
      \begin{atencao}
        Cobertura mede o que foi \textbf{executado}, não o que foi \textbf{conferido}: um teste sem \cod{expect} também ``cobre''.
      \end{atencao}
    \end{column}
    \begin{column}{0.52\textwidth}
      \scriptsize
      \begin{tabularx}{\linewidth}{@{}p{2.1cm} L@{}}
        \toprule
        \cabecalho{Critério} & \cabecalho{Exige, no exemplo} \\
        \midrule
        \textbf{Comandos} & cada linha rodar ao menos uma vez: 1 teste (30\,°C, ativo) basta \\
        \textbf{Decisões} & cada \cod{se} (\emph{branch}) ser verdadeiro e falso: 2 testes \\
        \textbf{Condições} & cada condição simples verdadeira e falsa: \cod{temp} e \cod{alertaAtivo} variando \\
        \textbf{Caminhos} & todas as rotas possíveis; cresce exponencialmente com laços e \cod{se}s \\
        \bottomrule
      \end{tabularx}
      \begin{conceito}[Hierarquia]
        Caminhos $\Rightarrow$ decisões $\Rightarrow$ comandos: quem satisfaz um critério mais forte satisfaz os mais fracos.
      \end{conceito}
    \end{column}
  \end{columns}
\end{frame}

\begin{frame}{Dublês de teste}
  \begin{columns}[T]
    \begin{column}{0.62\textwidth}
      \scriptsize
      \begin{tabularx}{\linewidth}{@{}l L L@{}}
        \toprule
        \cabecalho{Dublê} & \cabecalho{O que é} & \cabecalho{Exemplo no LabMonitor} \\
        \midrule
        \textbf{Dummy} & objeto só para preencher um parâmetro; nunca é usado & um \emph{logger} vazio exigido pelo construtor \\
        \textbf{Stub} & devolve respostas prontas & repositório que sempre devolve 3 bancadas \\
        \textbf{Spy} & stub que também registra as chamadas & conta quantas notificações foram enviadas \\
        \textbf{Mock} & objeto com \textbf{expectativas} de interação, verificadas no teste & ``\cod{reservar(7)} deve ser chamado uma vez'' \\
        \textbf{Fake} & implementação funcional simplificada & banco em memória, API falsa com dados fixos \\
        \bottomrule
      \end{tabularx}
    \end{column}
    \begin{column}{0.36\textwidth}
      \begin{definicao}[Dublê de teste]
        Objeto que substitui uma dependência real (rede, banco, sensor, relógio) durante o teste.
      \end{definicao}
      \begin{analogia}
        Como o dublê de cinema: entra nas cenas perigosas ou caras no lugar do ator.
      \end{analogia}
      \begin{historia}
        A classificação é de Gerard Meszaros, no livro \emph{xUnit Test Patterns} (2007).
      \end{historia}
    \end{column}
  \end{columns}
\end{frame}

\begin{frame}{Propriedades de um bom teste}
  \begin{columns}[T]
    \begin{column}{0.5\textwidth}
      \scriptsize
      \begin{tabularx}{\linewidth}{@{}l L@{}}
        \toprule
        \cabecalho{FIRST} & \cabecalho{Significado} \\
        \midrule
        \textbf{F}ast & rápido: roda em milissegundos, então roda sempre \\
        \textbf{I}ndependent & não depende da ordem nem de outros testes \\
        \textbf{R}epeatable & mesmo resultado em qualquer máquina e horário \\
        \textbf{S}elf-validating & passa ou falha sozinho, sem alguém olhar a saída \\
        \textbf{T}imely & escrito junto com o código (ou antes) \\
        \bottomrule
      \end{tabularx}
      \begin{conceito}[Estrutura]
        \textbf{AAA}: \emph{Arrange}, \emph{Act}, \emph{Assert}.\\
        \textbf{BDD}: \emph{Given}, \emph{When}, \emph{Then} (dado, quando, então).
      \end{conceito}
    \end{column}
    \begin{column}{0.48\textwidth}
      \begin{erro}[Teste instável (\emph{flaky})]
        Passa e falha sem o código mudar. Causas típicas: depender do relógio real, da rede, da ordem de execução, de esperas fixas (\emph{sleep}) ou de estado compartilhado entre testes.
      \end{erro}
      \begin{dica}[Teste também é código]
        Merece nomes claros, nada de duplicação e revisão. Um bom nome de teste descreve uma \textbf{regra}: ``reserva sobreposta é recusada''.
      \end{dica}
    \end{column}
  \end{columns}
\end{frame}

\begin{frame}[fragile]{TDD, BDD e regressão}
  \begin{columns}[T]
    \begin{column}{0.4\textwidth}
      \centering
      \begin{tikzpicture}
        \node[circle,fill=vermelhoerro,text=white,font=\scriptsize\bfseries,minimum size=1.35cm,align=center] (v) at (90:1.35) {Vermelho\\ \tiny teste falha};
        \node[circle,fill=verdeok,text=white,font=\scriptsize\bfseries,minimum size=1.35cm,align=center] (g) at (-30:1.35) {Verde\\ \tiny passa};
        \node[circle,fill=azulpetroleo,text=white,font=\scriptsize\bfseries,minimum size=1.35cm,align=center] (r) at (210:1.35) {Refatorar\\ \tiny melhora};
        \draw[setaambar] (v) to[bend left=25] (g);
        \draw[setaambar] (g) to[bend left=25] (r);
        \draw[setaambar] (r) to[bend left=25] (v);
      \end{tikzpicture}
      \begin{definicao}[TDD]
        \emph{Test-Driven Development} (Kent Beck): o teste é escrito \textbf{antes} do código e guia o projeto.
      \end{definicao}
    \end{column}
    \begin{column}{0.58\textwidth}
\begin{codigo}[title=BDD: cenário em linguagem natural (Gherkin)]{bash}
Funcionalidade: Reserva de bancada
  Cenário: horário já ocupado
    Dado que a bancada B7 está reservada das 14h às 16h
    Quando eu peço a B7 das 15h às 17h
    Então vejo "Horário indisponível"
\end{codigo}
      \begin{definicao}[BDD]
        \emph{Behavior-Driven Development}: descreve o comportamento em exemplos que clientes, testadores e devs entendem; os cenários viram testes.
      \end{definicao}
      \begin{definicao}[Teste de regressão]
        Rodar de novo os testes existentes após cada mudança, para garantir que nada que funcionava quebrou. Automatizar é o que torna isso viável.
      \end{definicao}
    \end{column}
  \end{columns}
\end{frame}

\begin{frame}{Desafios do teste em apps móveis}
  \begin{columns}[T]
    \begin{column}{0.5\textwidth}
      \scriptsize
      \begin{tabularx}{\linewidth}{@{}p{2.0cm} L@{}}
        \toprule
        \cabecalho{Desafio} & \cabecalho{O que pode dar errado} \\
        \midrule
        Fragmentação & milhares de modelos Android, versões de sistema, tamanhos e densidades de tela \\
        Rede & sem sinal, 3G lento, troca Wi-Fi/móvel no meio de uma requisição \\
        Interrupções & ligação, notificação, app enviado ao fundo e encerrado pelo sistema \\
        Hardware & câmera, GPS, Bluetooth e sensores diferentes ou ausentes \\
        Permissões & usuário nega, concede ``só desta vez'' ou revoga depois \\
        Recursos & bateria fraca, pouca memória, armazenamento cheio \\
        Pessoas & idioma, fonte grande, leitor de tela, modo escuro \\
        \bottomrule
      \end{tabularx}
    \end{column}
    \begin{column}{0.48\textwidth}
      \begin{conceito}[Estratégia]
        \begin{itemize}
          \item Automatizar a lógica em testes de unidade e widget (rápidos, sem aparelho).
          \item Definir uma \textbf{matriz de dispositivos} representativa (versões antigas e novas, tela pequena e grande).
          \item Usar emuladores no dia a dia e aparelhos reais ou \emph{device farms} antes de publicar.
          \item Simular rede ruim e interrupções.
        \end{itemize}
      \end{conceito}
      \begin{dica}
        As estatísticas do Android Studio e do console da loja mostram quais versões seus usuários realmente usam.
      \end{dica}
    \end{column}
  \end{columns}
\end{frame}

\begin{frame}{Qualidade é mais do que testar}
  \begin{columns}[T]
    \begin{column}{0.5\textwidth}
      \begin{definicao}[Modelo de qualidade ISO/IEC 25010]
        Organiza a qualidade de um produto de software em características como adequação funcional, eficiência de desempenho, compatibilidade, usabilidade, confiabilidade, segurança, manutenibilidade e portabilidade.
      \end{definicao}
      \begin{dica}
        A revisão de 2023 renomeou algumas características (usabilidade virou capacidade de interação) e incluiu \emph{safety}.
      \end{dica}
    \end{column}
    \begin{column}{0.48\textwidth}
      \begin{conceito}[Técnicas que se somam]
        \begin{itemize}
          \item \textbf{Estáticas} (sem executar): revisão de código, análise estática, \emph{lint}, tipagem forte.
          \item \textbf{Dinâmicas} (executando): testes automatizados e exploratórios.
          \item \textbf{Em produção}: monitoramento de falhas, métricas e \emph{feedback} dos usuários.
        \end{itemize}
      \end{conceito}
      \begin{destaque}
        Qualidade se \textbf{constrói} durante o desenvolvimento; testes confirmam que ela está lá.
      \end{destaque}
    \end{column}
  \end{columns}
\end{frame}
